\chapter{Smooth morphisms: generalities, differential properties}
\label{II}

% original p. 29
Cross-references to Expos\'e~\SourceRef{I}{I} are indicated by~I. One
recalls that rings are noetherian, and preschemes locally noetherian.

\section{Generalities}
\label{II.1}

Let $Y$ be a prescheme, let $t_1,\dots,t_n$ be
indeterminates, one sets
\begin{equation}
\label{eq:II.1.1}
Y[t_1,\dots,t_n]=Y\otimes_\ZZ \ZZ[t_1,\dots,t_n]\quoi.
\end{equation}
Thus $Y[t_1,\dots,t_n]$ is a $Y$-scheme, affine over~$Y$,
defined by the quasi-coherent sheaf of algebras
$\cal{O}_Y[t_1,\dots,t_n]$. The datum of a section of this
prescheme over~$Y$ is therefore equivalent to the
datum of~$n$ sections of~$\cal{O}_Y$ (corresponding to the images
of the~$t_i$ under the corresponding homomorphism). If $Y'$ is
over~$Y$, one has
\begin{equation}
\label{eq:II.1.2}
Y[t_1,\dots,t_n]\times_Y Y'=Y'[t_1,\dots,t_n]\quoi,
\end{equation}
(which implies that the datum of a $Y$-morphism from~$Y'$
to $Y[t_1,\dots,t_n]$ is equivalent to the datum of~$n$
sections of~$\cal{O}_{Y'}$), on the other hand one has
\begin{equation}
\label{eq:II.1.3}
\big(Y[t_1,\dots,t_n]\big)[t_{n+1},\dots,t_m]=Y[t_1,\dots,t_m]\quoi,
\end{equation}
by virtue of the analogous formula for polynomial rings
over~$\ZZ$. The formula~\eqref{eq:II.1.2} implies that
$Y[t_1,\dots,t_n]$ varies functorially with~$Y$.

$Y[t_1,\dots,t_n]$ is of finite type and flat over~$Y$.

\begin{definition}
\label{II.1.1}
Let $f\colon X \to Y$ be a morphism, making~$X$ into a
$Y$-prescheme. One says that $f$ is
\emph{smooth}\footnote{Old terminology: $f$ is \emph{simple}
at~$x$, or $x$ is a \emph{simple} point for~$f$. This
terminology led to confusion in various contexts
(simple algebras, simple groups) and had to be
abandoned.}
at $x \in X$, or that $X$ is \emph{smooth over~$Y$ at~$x$}, if there exists
an integer $n \geq 0$, an open neighborhood $U$ of~$x$, and an
\'etale $Y$-morphism from~$U$ to $Y[t_1,\dots,t_n]$. One says
that~$f$
(\resp $X$)
is \emph{smooth} if it is smooth at every point of~$X$. An
algebra~$B$ over a ring~$A$ is said to be smooth at a prime
ideal $\goth{p}$ of~$B$, if $\Spec(B)$ is smooth over~$\Spec(A)$
at
the point~$\goth{p}$;
% original p. 30
$B$ is said to be smooth over~$A$ if it is smooth over~$A$ at every prime
ideal~$\goth{p}$ of~$B$. Finally, a local homomorphism $A \to B$
of local rings is said to be smooth (or $B$ is said to be smooth over~$A$)\footnote{It is then better to say, as in EGA IV 18.6.1, that $B$
is ``\emph{essentially smooth}'' over~$A$.}
if $B$ is a localization of an algebra of finite type $B_1$ smooth
over~$A$.
\end{definition}

One notes that the notion of smoothness of~$X$ over~$Y$ is local on~$X$
and on~$Y$; if $X$ is smooth over~$Y$, it is locally of finite type
over~$Y$.

% The corrected source repeats 1.1. Preserve it with a distinct PDF anchor.
\setcounter{theorem}{0}
\begingroup
\renewcommand{\theHtheorem}{\theHsection.\arabic{theorem}.bis}
\begin{proposition}
\label{prop:II.1.1}
The set of points $x$ of~$X$ at which $f$ is smooth is open.
\end{proposition}
\endgroup

This is trivial from the definition.

\begin{corollary}
\label{II.1.2}
If $B$ is smooth over~$A$ at~$\goth{p}$, then it is smooth over~$A$
at~$\goth{q}$ for every prime ideal~$\goth{q}$ of~$B$ contained
in~$\goth{p}$.
\end{corollary}

\ref{prop:II.1.1} also implies that the last two
definitions~\ref{II.1.1} coincide in their common
domain of existence.

\begin{proposition}
\label{II.1.3}
\textup{(i)}~An \'etale morphism, in particular an open
immersion, an identity morphism, is smooth. \textup{(ii)}~An extension
of the base in a smooth morphism gives a smooth
morphism. \textup{(iii)}~The composite of two smooth morphisms is
smooth.
\end{proposition}

(i)~is trivial from the definition, more precisely one has:

\begin{corollary}
\label{II.1.4}
\'etale $=$ quasi-finite $+$ smooth.
\end{corollary}

(ii)~follows at once from the analogous fact for
\'etale morphisms (I~\SourceRef{I.4.6}{4.6}) and for the projections $Y[t_1,\dots,t_n]
\to Y$ (\cf \eqref{eq:II.1.2}). For (iii), this follows
formally from the fact that it is true separately for
``\'etale'' (I~\SourceRef{I.4.6}{4.6}) and for projections of the type
$Y[t_1,\dots,t_n]\to Y$
(\cf \eqref{eq:II.1.3}), and from the two facts cited for~(ii):
Suppose $Y$ is smooth over~$Z$ and $X$ is smooth over~$Y$, let us prove that $X$ is
smooth over~$Z$; one may suppose $Y$ \'etale over~$Z[t_1,\dots,t_n]$
and $X$ \'etale over~$Y[s_1,\dots,s_m]$, the first
hypothesis therefore implies that $Y[s_1,\dots,s_m]$ is \'etale over~$Z[t_1,\dots,t_n][s_1,\dots,s_m] = Z[t_1,\dots,s_m]$, hence $X$ is
\'etale over~$Z[t_1,\dots,s_m]$, qed.

\begin{remark}
\label{II.1.5}
The integer $n$ appearing in def\ptbl \ref{II.1.1} is well
determined, for one sees
% original p. 31
at once that it is the dimension of the local ring of~$x$ in its
fiber $f^{-1}\big(f(x)\big)$. One calls it the ``relative dimension''
of~$X$ over~$Y$. It behaves additively for the composition of
morphisms.
\end{remark}

\section{Some criteria of smoothness of a morphism}
\label{II.2}

\begin{theorem}
\label{II.2.1}
Let $f\colon X\to Y$ be a morphism locally of finite type, let $x\in
X$ and $y=f(x)$. For $f$ to be smooth at $x$, it is necessary and sufficient
that \textup{(a)} $f$ be flat at $x$, and \textup{(b)} $f^{-1}(y)$ be smooth over~$\kres(y)$
at $x$.
\end{theorem}

The composite of two flat morphisms being flat, and
$Y[t_{1},\dots,t_{n}]\to Y$ being a flat morphism, one sees that
smooth implies flat; taking~\ref{II.1.3}~(ii) into account this proves
necessity. Suppose (a) and (b) are satisfied, let $V$ be
an affine neighborhood of~$y$ with ring~$A$, $U$ an affine neighborhood of
$x$ over~$V$, with ring $B$. Taking $U$ small enough, one may
suppose by (b) that there exists a $\kres(y)$-morphism \emph{\'etale}
\[
g\colon U|f^{-1}(y)\to\Spec k[t_{1},\dots,t_{n}]\quad (k=\kres(y))
\]
defined by $n$ sections $g_{i}$ of the structure sheaf of
$U|f^{-1}(y)$. One easily sees that one may suppose that the
$g_{i}$ (which a priori are elements of
$B\otimes_{A}k=BS^{-1}$, where $S=A-\goth{p}$, $\goth{p}$ the prime
ideal of~$A$ corresponding to $y$) come from sections of
the structure sheaf of
$U$,
hence that $g$ is induced by a morphism,
still denoted~$g$
\[
g\colon U\to Y[t_{1},\dots,t_{n}]
\]
(up to multiplying the $g_{i}$ by one and the same
nonzero element of~$k$). Now
$U$
is flat over~$Y$ by (a), the same is true of
$Y[t_{1},\dots,t_{n}]$, on the other hand $g$ induces an
\'etale morphism between the fibers over~$y$, hence $g$ is \'etale at
$x$ by~(I~\SourceRef{I.5.8}{5.8}), qed.

\begin{corollary}
\label{II.2.2}
Let $S$ be a prescheme, $f\colon X\to Y$ an $S$-morphism of
finite type, $Y$ being of finite type and flat over~$S$, $x\in X$, $s$
the projection of~$x$ on~$S$. For $f$ to be smooth at $x$, it is necessary
and sufficient that $X$ be flat (or again: smooth) over~$S$ at $x$, and
that the morphism $f_{s}\colon X_{s}\to Y_{s}$ induced on the fibers of
$s$ be smooth at $x$.
\end{corollary}

Only sufficiency requires a proof, and follows from
the criterion~\ref{II.2.1}, together with the
flatness criterion~(I~\SourceRef{I.5.9}{5.9}).

To
% original p. 32
state the following result, ``recall'' that a morphism
$f\colon X\to Y$ locally of finite type is said to be
\emph{equidimensional}
at the point $x\in X$ if (setting $y=f(x)$) one can find an open
neighborhood $U$ of~$x$, every component of which dominates a component of~$Y$
such that, for every $y'\in Y$, the irreducible components of
$f^{-1}(y')\cap U$ all have one and the same dimension
independent of~$y'$. It is enough, moreover, in this condition to
take for $y'$ the generic points of the irreducible
components of~$Y$ passing through $y$, \emph{and} the point $y$. If
for example $X$ and $Y$ are integral and $f$ is dominant, the condition
means that the components of the $f^{-1}(y)$ passing through $x$ have ``the
right'' dimension, \ie the dimension of the generic fiber
(recall that they are always $\ge$ the dimension of the
generic fiber). If $f$ is equidimensional at $x$, the
dimension of its fiber at $x$ being $n$, and if $g\colon U\to
Y'=Y[t_{1},\dots,t_{n}]$ is a $Y$-morphism of a neighborhood~$U$ of
$x$, inducing a morphism on the fibers of~$y$ which is quasi-finite
at $x$ (or again, which amounts to the same, if $g$ is quasi-finite
at~$x$), then one shows that every irreducible component of~$U$
passing through $x$ dominates an irreducible component of~$Y'$. Moreover by virtue of the ``normalization lemma'', such a $g$
always exists (and conversely, if there exists a quasi-finite $Y$-morphism
$g$ of an open neighborhood $U$ of~$x$ into a $Y$-scheme
of the form $Y'=Y[t_{1},\dots,t_{n}]$, such that every component of~$U$ passing through $x$ dominates a component of~$Y'$, then $f$ is
equidimensional at~$x$). This being so:

\begin{proposition}
\label{II.2.3}
Let $f\colon X\to Y$ be a morphism locally of finite type, $x$ a
point of~$X$, $y=f(x)$, one assumes $\cal{O}_{y}$ normal. For $f$
to be smooth at $x$, it is necessary and sufficient that $f$ be
equidimensional at $x$, and that $f^{-1}(y)$ be smooth over~$\kres(y)$
at $x$.
\end{proposition}

One sees at once from the definition that a smooth morphism is
equidimensional (N.B. a flat morphism of finite type is not
necessarily equidimensional at $x$, even if its fiber
at $x$ is irreducible). Let us prove the converse. Since
$f^{-1}(y)$ is smooth over~$\kres(y)$ at $x$, one may suppose (replacing $X$ if need be by a suitable neighborhood of~$x$) that there exists
a $Y$-morphism
\[
g\colon X\to Y[t_{1},\dots,t_{n}]=Y'
\]
inducing an \'etale morphism on the fibers of~$y$, and a fortiori
quasi-finite at $x$. Hence
% original p. 33
$g$ is unramified, and ($f$ being equidimensional at
$x$) the irreducible components of~$X$ passing through $x$ each dominate
a component of~$Y'$, a fortiori the homomorphism
$\cal{O}_{y'}\to\cal{O}_{x}$ deduced from~$g$ (where $y'=g(x)$) is
\emph{injective}. This homomorphism is moreover unramified, and
$\cal{O}_{y'}$ is normal since it is a localization of the ring
$\cal{O}_{y}[t_{1},\dots,t_{n}]$, which is normal since
$\cal{O}_{y}$ is so. Hence the homomorphism $\cal{O}_{y'}\to\cal{O}_{x}$
is \'etale (I~\SourceRef{I.9.5}{9.5}~(ii)).

\begin{remarks}
\label{II.2.4}
The preceding statement still holds upon replacing
the hypothesis that $\cal{O}_{y}$ is normal by the weaker
hypothesis that $Y$ is \emph{geometrically unibranch} at $y$,
(\cf I~\SourceRef{I.11}{11}) --- since (I~\SourceRef{I.9.5}{9.5}) holds under this
hypothesis. Let us take this opportunity to note at the same time
that if the residue field of a local integral domain $A$ is
algebraically closed, then analytically integral
(\ie $\widehat{A}$ is an integral domain) implies geometrically
unibranch, the converse being moreover true in any
category of ``good rings'', more precisely in a
category of rings stable under the usual operations, and
in which the completion of a normal local ring is normal
(a condition fulfilled, by virtue of the ``analytic normality
theorem'' of Zariski, in the category of affine
algebras and their localizations)\footnote{\Cf EGA IV~7.8.}.

``Recall'' finally in the present context the following result,
due to Hironaka\footnote{\Cf EGA IV 5.12.10.} which sometimes allows
one to ensure that $f^{-1}(y)$ is a reduced scheme, \ie that
it is also what many algebraic geometers
would abusively regard as the fiber (without multiplicity)
of~$f$ above~$x$ (namely $f^{-1}(y)_{\red}$):
\end{remarks}

\begin{proposition}
\label{II.2.5}
Let $f\colon X\to Y$ be a dominant morphism of finite type of
reduced preschemes, $y$ a point of~$Y$ such that
$\cal{O}_{y}$ is regular. One assumes that all the components
of~$f^{-1}(y)$ are of multiplicity~$1$
(\cf the definition below), and that $f^{-1}(y)_{\red}$ is
normal. Then $f^{-1}(y)$ is reduced hence normal, $X$ is normal
at all the points of~$f^{-1}(y)$, finally $X$ is flat over~$Y$ at all
the points of~$f^{-1}(y)$.
\end{proposition}
One
% original p. 34
says that a component $Z$ of~$f^{-1}(y)$ is of
\emph{multiplicity}~$1$
if, $x$ denoting the generic point of~$Z$, one has (i)
$\dim\cal{O}_{x}=\dim\cal{O}_{y}$ (\ie $Z$ is not an ``excess
component'', that is, is not ``of too large
dimension''; (ii) the maximal ideal of~$\cal{O}_{x}$ is generated
by the maximal ideal of~$\cal{O}_{y}$ (which a priori, by virtue of
the choice of~$x$, generates an ideal of definition
of~$\cal{O}_{x}$).

Taking~\ref{II.2.3} or~\ref{II.2.1} into account one thus finds:

\begin{corollary}
\label{II.2.6}
Let $f\colon X\to Y$ be a dominant morphism of finite type of
reduced preschemes, $y$ a point of~$Y$ such that
$\cal{O}_{y}$ is regular. For $f$ to be smooth at the points of
$X$ above~$y$, it is necessary and sufficient that the components of
$f^{-1}(y)$ be of multiplicities $1$, and that $f^{-1}(y)_{\red}$
be smooth over~$\kres(y)$.
\end{corollary}

This situation was considered especially in the past
when $Y$ was the spectrum of a discrete valuation ring
$A$, and was commonly designated by phrases
such as: ``if the reduction of~$X$ with respect to the given
valuation is nice''... Moreover, $X$ then designated a
closed subscheme (if one may say so) of a $\PP_{K}^{n}$ ($K$
being the field of fractions of~$A$) and for want of an
adequate language, the more intrinsic role of an object ``defined
over~$A$'' (and not merely over~$K$) hardly appeared.

\section{Permanence properties}
\label{II.3}

\begin{proposition}
\label{II.3.1}
Let $f\colon X\to Y$ be a morphism, let $x\in X$ and
$y=f(x)$. Suppose $f$ is smooth at~$x$. For $\cal{O}_{x}$ to be
reduced (\resp regular, \resp normal) it is necessary and sufficient that
$\cal{O}_{y}$ be so.
\end{proposition}

This statement is indeed known when $X$ is of the form
$Y[t_{1},\dots,t_{n}]$, and it is proved in (I, \no \SourceRef{I.9}{9})
for an \'etale morphism; the general case is deduced from this
at once thanks to definition~\ref{II.1.1}.

We do not spell out here the other permanence
properties, already resulting from flatness alone, or from
the fact that $X$ is locally quasi-finite and flat over a
$Y$-prescheme of the form $Y[t_{1},\dots,t_{n}]$ (or, as
% original p. 35
we shall say, that $X$ is
Cohen--Macaulay
over~$Y$). Let us only note that from this last fact
it follows that
\begin{equation}
\label{eq:II.3.1}
{\dim\cal{O}_{x}=\dim\cal{O}_{y}+n-d,\quad
\prof\cal{O}_{x}=\prof\cal{O}_{y}+n-d},
\end{equation}
where $n$ is the dimension of the fiber of~$f$ at $x$, and $d$ the
transcendence degree of~$\kres(x)$ over~$\kres(y)$, whence (setting
$\coprof=\dim-\prof$)\footnote{For these formulas, \cf EGA~IV 6.1
and~6.3.}
\begin{equation}
\label{eq:II.3.2}
{\coprof\cal{O}_{x}=\coprof\cal{O}_{y}}.
\end{equation}
It follows for example that $\cal{O}_{x}$ is Cohen--Macaulay
(\resp without embedded components) if and only if the same is
true of~$\cal{O}_{y}$.
